\chapter{Limitações e discussão}

\section{Limitações da linguagem}

A MiniLang continua restrita ao nível lexical. O reconhecimento de um identificador não determina se ele foi declarado ou qual seu papel semântico. O reconhecimento de um número verifica somente seu formato textual. Horários são validados segundo as faixas 00--23 e 00--59, mas não possuem interpretação contextual.

Strings não possuem escapes na primeira versão. Isso simplifica a construção do AFN$\varepsilon$, mas limita a representação de conteúdos que exijam aspas ou barra invertida internas.

\section{Limitações do lexer}

O MiniLexer não constrói árvore sintática nem interpreta a estrutura do programa. O arquivo pode ser lexicalmente válido sem ser um programa sintaticamente válido em uma linguagem de programação completa.

\section{Melhorias futuras}

Como evolução do projeto, podem ser considerados: introdução de escapes em strings com uma nova ER e correspondente AFN$\varepsilon$; expansão controlada do vocabulário da MiniLang; criação de uma camada de análise sintática; e ampliação das baterias de testes para alfabetos maiores. Essas extensões não fazem parte da versão entregue e, portanto, não foram utilizadas para produzir os resultados atuais.

\section{Decisões de escopo}

Optou-se por não acrescentar funcionalidades apenas para tornar as ERs artificialmente mais complexas. As seis ERs existentes já superam o mínimo exigido e são diretamente justificadas pelo vocabulário lexical da aplicação.
